% sec:gr-exact-solutions — Exact Solutions: An Overview
\section{Exact Solutions: An Overview}
\label{sec:gr-exact-solutions}

The field equations tell us how geometry responds to matter, but to
integrate geodesics we need an explicit metric at every spacetime
point.  Finding exact solutions to the vacuum equations $\ricci = 0$
(\cref{eq:vacuum-efe}) is notoriously difficult --- the nonlinear
coupling means that even elementary symmetry assumptions lead to
lengthy calculations --- yet the known solutions include precisely the
spacetimes that generate the most dramatic images: non-rotating and
rotating black holes, and traversable wormholes.  Three exact solutions
are implemented in the codebase, each corresponding to a physically
distinct geometry, and they provide the input for everything the ray
tracer computes.
\coderef{Spacetime}{Class}

\paragraph{The Schwarzschild metric.}

The simplest black hole solution assumes no rotation and no charge.
Birkhoff's theorem (\cref{sec:schw-derivation}) guarantees that
spherical symmetry alone, combined with the vacuum equations,
determines the metric uniquely up to the mass parameter~$M$.  In
Boyer--Lindquist coordinates $(t, r, \theta, \varphi)$, the line
element is
%
\begin{equation}
  \d s^2 = -\!\left(1 - \frac{2M}{r}\right)\d t^2
           + \left(1 - \frac{2M}{r}\right)^{-1}\!\d r^2
           + r^2\!\left(\d\theta^2 + \sin^2\!\theta\;\d\varphi^2\right) .
  \label{eq:schwarzschild-metric}
\end{equation}
%
The time--time component $g_{00} = -(1 - 2M/r)$ is the gravitational
potential identified in the weak-field analysis (\cref{eq:g00-potential});
the exact redshift formula (\cref{eq:redshift-schwarzschild}) follows
directly from it.  Three radii organise the physics: the event
horizon at $r = 2M$, where photons are infinitely redshifted; the
photon sphere at $r = 3M$, where unstable circular photon orbits
produce the bright ring visible in resolved black hole images; and the
innermost stable circular orbit (ISCO) at $r = 6M$, the inner edge
of a thin accretion disc (\cref{sec:schw-orbits}).  For a
$10\,M_\odot$ stellar-mass black hole, these correspond to roughly
$30\;\text{km}$, $44\;\text{km}$, and $89\;\text{km}$ --- city-scale
distances for an object containing ten solar masses.
\histnote{Schwarzschild found this solution in December 1915, just
weeks after Einstein published the field equations --- while serving
on the Eastern Front in World War~I.  He communicated the result
to Einstein by letter; it was the first exact solution.}

The divergence of $g_{rr}$ at $r = 2M$ is a coordinate artefact, not
a physical singularity: the Kretschmann scalar
$R_{\alpha\beta\gamma\delta}R^{\alpha\beta\gamma\delta} = 48M^2/r^6$
is finite there and diverges only at $r = 0$.  But for the ray tracer,
coordinates that break at the horizon are unacceptable --- photons
must be tracked smoothly through $r = 2M$.
\xrefnote{The full treatment of coordinate singularities and
their resolution appears in \cref{sec:schw-coordinates}.}

\paragraph{The Kerr--Schild form.}

The coordinate singularity at $r = 2M$ demands a metric form that
remains finite across the horizon.  Both the Schwarzschild and Kerr
spacetimes admit such a form --- a decomposition in which the full
metric is a flat background plus a single tensor correction:
%
\begin{equation}
  \metric = \eta_{\mu\nu} + f\,l_\mu\,l_\nu \,,
  \label{eq:kerr-schild-form}
\end{equation}
%
where $f$ is a scalar function and $l_\mu$ is a null covector with
respect to both $\eta_{\mu\nu}$ and $\metric$ --- a property unique
to the Kerr--Schild class.  Because $l_\mu$ is null with respect to
$\eta_{\mu\nu}$, the cross-terms that would normally appear in a
matrix inversion vanish, and the inverse metric takes a remarkably
simple form: $\inversemetric = \eta^{\mu\nu} - f\,l^\mu\,l^\nu$.
For the Schwarzschild spacetime in Cartesian coordinates
$(t, x, y, z)$, the Kerr--Schild data are $f = 2M/r$ and
$l_\mu = (1,\; x/r,\; y/r,\; z/r)$ with $r = \sqrt{x^2 + y^2 + z^2}$.
Every component of $\metric$ remains finite at $r = 2M$, so geodesics
can be integrated without special treatment as they cross the horizon.
The codebase uses Kerr--Schild coordinates as its default for both
Schwarzschild and Kerr spacetimes (\cref{sec:schw-kerr-schild}).
\coderef{Spacetime}{SchwarzschildKS}

\paragraph{The Kerr metric.}

Astrophysical black holes rotate.  The Kerr solution adds a single
parameter --- the specific angular momentum $a = J/M$ --- and breaks
spherical symmetry down to axisymmetry.  The Boyer--Lindquist line
element (derived in \cref{ch:kerr}) is governed by two functions:
%
\begin{equation}
  \Sigma = r^2 + a^2 \cos^2\!\theta \,, \qquad
  \Delta = r^2 - 2Mr + a^2 \,,
  \label{eq:kerr-sigma-delta}
\end{equation}
%
The function $\Sigma$ encodes the oblate-spheroid geometry
introduced by the spin; $\Delta$ is a quadratic in $r$ whose roots
are the event horizons.  Setting $\Delta = 0$ yields two horizons:
the outer event horizon at $r_+ = M + \sqrt{M^2 - a^2}$ and the
inner (Cauchy) horizon at $r_- = M - \sqrt{M^2 - a^2}$.  Both exist
only for $\abs{a} \leq M$ --- the bound separating black holes from
naked singularities.  The limiting case $\abs{a} = M$ is an extremal
Kerr black hole, whose two horizons merge at $r = M$.

Rotation introduces physics absent from the Schwarzschild case.
The off-diagonal component $g_{t\varphi}$ couples time and azimuth,
dragging spacetime along with the black hole --- an effect called
\emph{frame dragging}.  Inside the \emph{ergosphere}, frame dragging
is so strong that no observer can remain stationary relative to
infinity.  For the ray tracer, this produces the characteristic
brightness asymmetry of a rotating black hole: co-rotating photons
are blueshifted while counter-rotating photons are redshifted.

The Kerr metric also admits the Kerr--Schild form
(\cref{eq:kerr-schild-form}), with
$f = 2Mr^3/(r^4 + a^2 z^2)$ --- where $r$ is now the Boyer--Lindquist
radius, an implicit function of the Cartesian coordinates that reduces
to the Euclidean radius when $a = 0$ --- and a null covector $l_\mu$
adapted to the spin.  Setting $a = 0$ recovers the Schwarzschild case
exactly.
\coderef{Spacetime}{KerrKS}

\paragraph{The Morris--Thorne wormhole.}

Not every exact solution describes a black hole.  The Morris--Thorne
metric is a static, spherically symmetric spacetime with a
\emph{throat} --- a minimal-area sphere connecting two asymptotically
flat regions \cite{MorrisThorne1988}.  In terms of the proper radial
distance $\ell$ from the throat, the line element is
%
\begin{equation}
  \d s^2 = -\d t^2 + \d\ell^2
           + (b_0^2 + \ell^2)\!\left(\d\theta^2
             + \sin^2\!\theta\;\d\varphi^2\right) ,
  \label{eq:morris-thorne-metric}
\end{equation}
%
where $b_0$ is the throat radius.  The circumferential radius
$r(\ell) = \sqrt{b_0^2 + \ell^2}$ reaches its minimum $b_0$ at the
throat ($\ell = 0$) but never vanishes: there is no singularity and
no horizon.  Geodesics pass freely through the throat from one
asymptotic region to the other.

The general Morris--Thorne family includes a redshift function
$\Phi(\ell)$ multiplying $g_{tt}$; the codebase implements the
zero-redshift case $\Phi = 0$, giving a globally constant
$g_{tt} = -1$ and no gravitational redshift anywhere in the
spacetime.  This simplicity has a cost: sustaining the throat requires
matter that violates the null energy condition (\cref{eq:nec}), as
discussed in \cref{sec:gr-stress-energy}.  In the codebase, the
Morris--Thorne wormhole serves as a visually striking test case ---
the double sky visible through the throat and the lensing ring around
its rim provide a stringent check on the geodesic integrator's
handling of nontrivial topology.
\warnnote{No known classical or quantum matter satisfies the
NEC violation required to sustain a wormhole throat.  The
Morris--Thorne spacetime is a valid solution to the Einstein
equations, but the matter source it demands has no known physical
realisation --- it is a mathematical existence proof, not a
prediction.}
\coderef{Spacetime}{MorrisThorne}

\begin{figure}[ht]
  \centering
  % embedding-diagrams.tex — Side-by-side embedding diagrams:
%   (a) Schwarzschild black hole, (b) Morris-Thorne wormhole
% Inputable TikZ fragment for fig:embedding-diagrams
% Requires: tikz with calc, arrows.meta (loaded by ehbook.cls)
%
% Each panel shows the 2D cross-section (r, z) of the equatorial
% spatial slice embedded into Euclidean 3-space.
%
% Schwarzschild:  z(r) = ± 2 sqrt(2M(r - 2M)),  M = 1,  r in [2, 8]
% Morris-Thorne:  r(ell) = sqrt(b0^2 + ell^2),   b0 = 2, ell in [-7.5, 7.5]
%
\begin{tikzpicture}[
    >={Stealth[length=5pt]},
    font=\footnotesize,
  ]

  % ── Scale parameters ──────────────────────────────────────────────
  \def\rscale{0.45}      % cm per unit of circumferential radius r
  \def\zscale{0.33}      % cm per unit of embedding height z or ell
  \def\panelsep{4.8}     % horizontal offset between panel origins (cm)

  % ====================================================================
  %  Panel (a): Schwarzschild embedding
  % ====================================================================
  \begin{scope}

    % ── Axes ──────────────────────────────────────────────────────────
    \draw[->, EHMarginGray, thin]
      (0, 0) -- ({8.5*\rscale}, 0)
      node[right, font=\scriptsize, text=EHMarginGray] {$r$};
    \draw[->, EHMarginGray, thin]
      (0, {-2.8*\zscale}) -- (0, {7.5*\zscale})
      node[above, font=\scriptsize, text=EHMarginGray] {$z$};

    % ── Upper branch: z(r) = +2 sqrt(2(r - 2)) ──────────────────────
    \draw[EHAccretionGold, thick] plot[smooth, tension=0.55] coordinates {
      ({2.00*\rscale}, {0.0000*\zscale})   % r=2.00  z=0.000 (horizon)
      ({2.02*\rscale}, {0.4000*\zscale})   % r=2.02  z=0.400
      ({2.05*\rscale}, {0.6325*\zscale})   % r=2.05  z=0.632
      ({2.10*\rscale}, {0.8944*\zscale})   % r=2.10  z=0.894
      ({2.20*\rscale}, {1.2649*\zscale})   % r=2.20  z=1.265
      ({2.40*\rscale}, {1.7889*\zscale})   % r=2.40  z=1.789
      ({2.70*\rscale}, {2.3664*\zscale})   % r=2.70  z=2.366
      ({3.00*\rscale}, {2.8284*\zscale})   % r=3.00  z=2.828
      ({3.50*\rscale}, {3.4641*\zscale})   % r=3.50  z=3.464
      ({4.00*\rscale}, {4.0000*\zscale})   % r=4.00  z=4.000
      ({4.50*\rscale}, {4.4721*\zscale})   % r=4.50  z=4.472
      ({5.00*\rscale}, {4.8990*\zscale})   % r=5.00  z=4.899
      ({5.50*\rscale}, {5.2915*\zscale})   % r=5.50  z=5.292
      ({6.00*\rscale}, {5.6569*\zscale})   % r=6.00  z=5.657
      ({6.50*\rscale}, {6.0000*\zscale})   % r=6.50  z=6.000
      ({7.00*\rscale}, {6.3246*\zscale})   % r=7.00  z=6.325
      ({7.50*\rscale}, {6.6332*\zscale})   % r=7.50  z=6.633
      ({8.00*\rscale}, {6.9282*\zscale})   % r=8.00  z=6.928
    };

    % ── Lower branch: z(r) = -2 sqrt(2(r - 2)) ──────────────────────
    \draw[EHAccretionGold, thick] plot[smooth, tension=0.55] coordinates {
      ({2.00*\rscale}, {-0.0000*\zscale})  % r=2.00  z=0.000 (horizon)
      ({2.02*\rscale}, {-0.4000*\zscale})  % r=2.02  z=-0.400
      ({2.05*\rscale}, {-0.6325*\zscale})  % r=2.05  z=-0.632
      ({2.10*\rscale}, {-0.8944*\zscale})  % r=2.10  z=-0.894
      ({2.20*\rscale}, {-1.2649*\zscale})  % r=2.20  z=-1.265
      ({2.40*\rscale}, {-1.7889*\zscale})  % r=2.40  z=-1.789
      ({2.70*\rscale}, {-2.3664*\zscale})  % r=2.70  z=-2.366
    };

    % ── Horizon marker ────────────────────────────────────────────────
    \draw[EHHorizonCrimson, dashed, thin]
      ({2*\rscale}, {-2.5*\zscale}) -- ({2*\rscale}, {1.5*\zscale});
    \fill[EHHorizonCrimson] ({2*\rscale}, 0) circle (1.5pt);
    \node[font=\scriptsize, text=EHHorizonCrimson, anchor=south east]
      at ({2*\rscale - 0.08}, {1.8*\zscale}) {$r\!=\!2M$};

    % ── Singularity arrow ─────────────────────────────────────────────
    \draw[->, EHHorizonCrimson, thick, shorten >=1pt]
      ({1.6*\rscale}, 0) -- ({0.2*\rscale}, 0);
    \node[font=\scriptsize, text=EHHorizonCrimson, anchor=east]
      at ({0.15*\rscale}, 0) {$r\!=\!0$};

    % ── Panel label ───────────────────────────────────────────────────
    \node[font=\small, anchor=north] at ({4.5*\rscale}, {-3.4*\zscale})
      {(a) Schwarzschild};

  \end{scope}

  % ====================================================================
  %  Panel (b): Morris-Thorne embedding
  % ====================================================================
  \begin{scope}[xshift=\panelsep cm]

    % ── Axes ──────────────────────────────────────────────────────────
    \draw[->, EHMarginGray, thin]
      (0, 0) -- ({8.5*\rscale}, 0)
      node[right, font=\scriptsize, text=EHMarginGray] {$r$};
    \draw[->, EHMarginGray, thin]
      (0, {-8.0*\zscale}) -- (0, {8.0*\zscale})
      node[above, font=\scriptsize, text=EHMarginGray] {$\ell$};

    % ── Embedding profile: (r(ell), ell) ─────────────────────────────
    %    r(ell) = sqrt(b0^2 + ell^2),  b0 = 2
    \draw[EHAccretionGold, thick] plot[smooth, tension=0.55] coordinates {
      ({7.7621*\rscale}, {-7.50*\zscale})  % ell=-7.5  r=7.762
      ({6.8007*\rscale}, {-6.50*\zscale})  % ell=-6.5  r=6.801
      ({5.8523*\rscale}, {-5.50*\zscale})  % ell=-5.5  r=5.852
      ({4.9244*\rscale}, {-4.50*\zscale})  % ell=-4.5  r=4.924
      ({4.0311*\rscale}, {-3.50*\zscale})  % ell=-3.5  r=4.031
      ({3.2016*\rscale}, {-2.50*\zscale})  % ell=-2.5  r=3.202
      ({2.5000*\rscale}, {-1.50*\zscale})  % ell=-1.5  r=2.500
      ({2.2361*\rscale}, {-1.00*\zscale})  % ell=-1.0  r=2.236
      ({2.0616*\rscale}, {-0.50*\zscale})  % ell=-0.5  r=2.062
      ({2.0000*\rscale}, { 0.00*\zscale})  % ell= 0.0  r=2.000 (throat)
      ({2.0616*\rscale}, { 0.50*\zscale})  % ell=+0.5  r=2.062
      ({2.2361*\rscale}, { 1.00*\zscale})  % ell=+1.0  r=2.236
      ({2.5000*\rscale}, { 1.50*\zscale})  % ell=+1.5  r=2.500
      ({3.2016*\rscale}, { 2.50*\zscale})  % ell=+2.5  r=3.202
      ({4.0311*\rscale}, { 3.50*\zscale})  % ell=+3.5  r=4.031
      ({4.9244*\rscale}, { 4.50*\zscale})  % ell=+4.5  r=4.924
      ({5.8523*\rscale}, { 5.50*\zscale})  % ell=+5.5  r=5.852
      ({6.8007*\rscale}, { 6.50*\zscale})  % ell=+6.5  r=6.801
      ({7.7621*\rscale}, { 7.50*\zscale})  % ell=+7.5  r=7.762
    };

    % ── Throat marker ─────────────────────────────────────────────────
    \draw[EHAccretionGold!80!black, dashed, thin]
      ({2*\rscale}, {-1.2*\zscale}) -- ({2*\rscale}, {1.2*\zscale});
    \fill[EHAccretionGold!80!black] ({2*\rscale}, 0) circle (1.5pt);
    \node[font=\scriptsize, text=EHAccretionGold!80!black, anchor=south east]
      at ({2*\rscale - 0.08}, {1.5*\zscale}) {$r\!=\!b_0$};

    % ── Universe labels ───────────────────────────────────────────────
    \node[font=\scriptsize, text=EHJetBlue, anchor=west]
      at ({8.0*\rscale}, {7.3*\zscale}) {universe I};
    \node[font=\scriptsize, text=EHJetBlue, anchor=west]
      at ({8.0*\rscale}, {-7.3*\zscale}) {universe II};

    % ── Panel label ───────────────────────────────────────────────────
    \node[font=\small, anchor=north] at ({4.5*\rscale}, {-8.6*\zscale})
      {(b) Morris--Thorne};

  \end{scope}

\end{tikzpicture}

  \caption[Embedding diagrams for Schwarzschild and Morris-Thorne spacetimes]{%
    Embedding diagrams for the Schwarzschild black hole (left) and the
    Morris--Thorne wormhole (right).  Each curve shows the cross-section
    of the equatorial ($\theta = \pi/2$) spatial slice embedded into
    Euclidean three-space.  The Schwarzschild geometry terminates at the
    event horizon $r = 2M$, where the funnel pinches off; the curvature
    singularity at $r = 0$ lies beyond.  The Morris--Thorne geometry has
    a smooth throat of finite circumference $2\pi b_0$ that connects two
    asymptotically flat regions.  Sustaining this throat geometry requires
    exotic matter that violates the null energy condition.}
  \label{fig:embedding-diagrams}
\end{figure}

\paragraph{Codebase architecture.}

The three spacetimes share a common interface: the
\texttt{Spacetime} type class.  Each instance provides the covariant
metric, its inverse, and the Christoffel symbols at any point:
%
\begin{lstlisting}[language=Haskell]
class Spacetime s where
  metric4      :: s -> Point s -> Sym4 'Covariant
  invMetric4   :: s -> Point s -> Sym4 'Contravariant
  christoffel4 :: s -> Point s -> Christoffel4
\end{lstlisting}
%
The type-level variance tags (\texttt{'Covariant},
\texttt{'Contravariant}) prevent accidental mixing of index positions
at compile time.  Because every spacetime exposes the same three
functions, the geodesic integrator developed in the following chapters
is entirely spacetime-agnostic: switching from Schwarzschild to Kerr
requires changing only the spacetime instance, not the integration
code.
\implnote{The Christoffel symbols are computed via automatic
differentiation by default.  Each spacetime provides a polymorphic
metric function of type \texttt{Floating~a~=>~[a]~->~[a]} that works
with both \texttt{Double} and AD types, eliminating hand-coded
derivatives entirely.  The AD strategy is developed in
\cref{ch:autodiff}.}

The exact solutions stated above are the inputs to the ray tracer;
the remaining chapters in this part derive each from first principles,
develop the geodesic machinery that turns $\metric$ and
$\christoffel{\lambda}{\mu}{\nu}$ into photon trajectories, and build
the spectral calculations that assign a colour and brightness to each
ray.
